\subsection*{1. Understanding the Given Data}
We are given a market with one riskless asset and two risky assets. 
The expected returns are given by the vector $R$, and the covariance matrix of the risky assets is $\Sigma$:
\[
R = \begin{pmatrix} 0.2 \\ 0.12 \end{pmatrix}, \quad 
\Sigma = \begin{pmatrix} 0.09 & 0.0045 \\ 0.0045 & 0.0225 \end{pmatrix}
\]
The riskless asset return is $R^0 = 0.08$. The problem states we are only investing in the risky assets, meaning the weight of the riskless asset is $w_0 = 0$. We also require a ``long position'' only, meaning our weights must be non-negative ($w_1, w_2 \ge 0$). 

Since we must fully invest our capital across the two risky assets, our weights must sum to 1:
\[
w_1 + w_2 = 1
\]

\subsection*{2. Calculating the Portfolio Variance}
The variance of a portfolio consisting of two assets, defined by the weight vector $w = (w_1, w_2)^T$, is given by $V(w) = w^T \Sigma w$. 

Let's expand this matrix multiplication step-by-step:
\[
V(w) = \begin{pmatrix} w_1 & w_2 \end{pmatrix} 
\begin{pmatrix} 0.09 & 0.0045 \\ 0.0045 & 0.0225 \end{pmatrix} 
\begin{pmatrix} w_1 \\ w_2 \end{pmatrix}
\]
First, multiply the row vector by the matrix:
\[
\begin{pmatrix} w_1 & w_2 \end{pmatrix} \Sigma = \begin{pmatrix} 0.09w_1 + 0.0045w_2 & 0.0045w_1 + 0.0225w_2 \end{pmatrix}
\]
Now, multiply this resulting row vector by the column vector $w$:
\[
V(w) = (0.09w_1 + 0.0045w_2)w_1 + (0.0045w_1 + 0.0225w_2)w_2
\]
\[
V(w) = 0.09w_1^2 + 0.0045w_1w_2 + 0.0045w_1w_2 + 0.0225w_2^2
\]
Combining the middle terms yields the portfolio variance equation:
\[
V(w) = 0.09w_1^2 + 0.009w_1w_2 + 0.0225w_2^2
\]

\subsection*{3. Formulating the Lagrangian}
We want to minimize the variance $V(w)$ subject to the constraint $w_1 + w_2 = 1$, which can be rewritten as $w_1 + w_2 - 1 = 0$. 

To do this, we introduce a Lagrange multiplier, $\lambda$, to combine our objective function and our constraint:
\[
L(w, \lambda) = V(w) + \lambda(w_1 + w_2 - 1)
\]
\[
L(w_1, w_2, \lambda) = 0.09w_1^2 + 0.009w_1w_2 + 0.0225w_2^2 + \lambda(w_1 + w_2 - 1)
\]

\subsection*{4. Solving the First-Order Conditions}
To find the minimum, we take the partial derivative of the Lagrangian with respect to $w_1, w_2,$ and $\lambda$, and set them equal to zero.

\vspace{0.5em}
\noindent \textbf{Derivative with respect to $w_1$:}
\[
\frac{\partial L}{\partial w_1} = 2(0.09)w_1 + 0.009w_2 + \lambda = 0.18w_1 + 0.009w_2 + \lambda = 0
\]

\noindent \textbf{Derivative with respect to $w_2$:}
\[
\frac{\partial L}{\partial w_2} = 0.009w_1 + 2(0.0225)w_2 + \lambda = 0.009w_1 + 0.045w_2 + \lambda = 0
\]

\noindent \textbf{Derivative with respect to $\lambda$:}
\[
\frac{\partial L}{\partial \lambda} = w_1 + w_2 - 1 = 0 \implies w_2 = 1 - w_1
\]

We now have a system of linear equations. Since both the $w_1$ and $w_2$ derivatives equal $-\lambda$, we can set them equal to each other:
\[
0.18w_1 + 0.009w_2 = 0.009w_1 + 0.045w_2
\]
Now, substitute $w_2 = 1 - w_1$ into this equation:
\[
0.18w_1 + 0.009(1 - w_1) = 0.009w_1 + 0.045(1 - w_1)
\]
\[
0.18w_1 + 0.009 - 0.009w_1 = 0.009w_1 + 0.045 - 0.045w_1
\]
\[
0.171w_1 + 0.009 = -0.036w_1 + 0.045
\]
Add $0.036w_1$ to both sides, and subtract $0.009$ from both sides:
\[
0.207w_1 = 0.036
\]
\[
\tilde{w}_1 = \frac{0.036}{0.207} \approx 0.1739
\]
Using $w_2 = 1 - w_1$, we find:
\[
\tilde{w}_2 = 1 - 0.1739 = 0.8261
\]
Since both weights are positive ($\tilde{w}_1, \tilde{w}_2 > 0$), this solution naturally satisfies the ``long only'' condition.

\subsection*{5. Expected Value \& Variance of the Return}
Now we calculate the expected return $r(V)$ and variance $Var(V)$ of this optimized portfolio.

\noindent \textbf{Expected Return:}
\[
r(V) = \tilde{w}^T R = (0.1739)(0.2) + (0.8261)(0.12) \approx 0.03478 + 0.099132 = 0.1339
\]
Note that this expected return ($13.39\%$) is greater than the risk-free rate $R^0 = 0.08$.

\noindent \textbf{Variance:}
Using our variance formula $V(w) = 0.09w_1^2 + 0.009w_1w_2 + 0.0225w_2^2$:
\[
Var(V) = 0.09(0.1739)^2 + 0.009(0.1739)(0.8261) + 0.0225(0.8261)^2
\]
\[
Var(V) \approx 0.00272 + 0.00129 + 0.01535 = 0.01936
\]
As expected, the minimum variance is strictly positive ($>0$), meaning uncertainty cannot be completely eliminated.